\documentclass[12pt]{article}
\usepackage[a4paper, margin=1in]{geometry}
\usepackage{amsmath, amssymb, minted, enumitem, cmbright, hyperref}
\renewcommand{\thesection}{\Alph{section}}
\renewcommand{\thesubsection}{\thesection.\arabic{subsection}}
\renewcommand{\t}{\texttt}
\begin{document}
\section{MCQs}
The answers are \textbf{bcd}...\textbf{aeb}.
\begin{enumerate}
\item List slicing is $O(n)$ in Python and hence the time complexity is $n+n/2+n/4+\cdots\in O(n)$. So the correct answer is b). $\checkmark$ (The correct implementation with $O(\log n)$ works with indices instead.)
\item The loop runs in $O(\log n)$ time, each iteration taking $O(n)$. So the correct answer is c). $\checkmark$
\item Creating $b$ takes $O\left(n^2\right)$ time which is the dominating term. So the correct answer is d). $\checkmark$

$\cdots$
\setcounter{enumi}{22}
\item It doesn't affect anything - the correct answer is a). $\checkmark$
\item By definition the correct answer is e). $\checkmark$
\item By definition the correct answer is b). $\checkmark$ (Surely it doesn't make sense...)
\end{enumerate}
All other questions are redacted.
\section{Simpler Questions}
\subsection{Feedback on IT5001}
(Omitted)
\subsection{Real-Life Application of DSA}
(Omitted)
\subsection{Triple-Ended Queue}
Maintain two deques, say $L$ and $R$, so that we always have $|L|=|R|$ when $|L|+|R|$ is even and $|L|=|R|+1$ when $|L|+|R|$ is odd, where $|L|$ and $|R|$ denote the size of $L$ and $R$ respectively.

Then we have, for a teque $T$,
\begin{itemize}
\item \t{T.append(x)} $\equiv$ \t{R.append(x)},
\item \t{T.appendleft(x)} $\equiv$ \t{L.appendleft(x)},
\item \t{T.insertmiddle(x)} $\equiv$ \t{R.appendleft(x)},
\item \t{T.printteque()} $\equiv$ \t{L.printdeque()} $+$ \t{R.printdeque()}.
\end{itemize}

After each step, we `rebalance' $T$ as needed to maintain the invariant: move the last element of $L$ to the front of $R$ if $L$ is too big, and vice versa.

A Python code is given below.
\begin{minted}{python}
from collections import deque

class Teque:
    def __init__(self):
        self.L = deque()
        self.R = deque()

    def _rebalance(self):
        if len(self.L) > len(self.R):
            self.R.appendleft(self.L.pop())
        elif len(self.R) > len(self.L) + 1:
            self.L.append(self.R.popleft())

    def append(self, x):
        self.R.append(x)
        self._rebalance()

    def appendleft(self, x):
        self.L.appendleft(x)
        self._rebalance()

    def insertmiddle(self, x):
        self.R.appendleft(x)
        self._rebalance()

    def printteque(self):
        for x in self.L:
            print(x, end=' ')
        for x in self.R:
            print(x, end=' ')
        print()
\end{minted}
\subsection{Stability of Binary (Max) Heap With Duplicates}
\subsubsection{Show Your Understanding!}
No. You can get $7_3,7_1,7_2$ for example. (VA OQ 1)
\subsubsection{Make Heap Sort a Stable Sort!}
Add an index for each duplicate element. For instance, $7_1,7_2,7_3$ could be inserted as pairs $(7,-1),(7,-2),(7,-3)$, etc., as it is a max heap. Then the order could be maintained.
\section{Applications}
\subsection{Graph Traversal}
\subsubsection{Check Your Understanding}
\begin{enumerate}
\item Tree
\item No. (...\textit{start from any source vertex}...)
\end{enumerate}
\subsubsection{Subtask 1: $k\ge2$}
Then we can go back and forth on edges and therefore traverse the entire tree. So the answer is simply the sum of all weights, $O(N)$.
\subsubsection{Subtask 2: $k=1$ and you must start from vertex $1$}
Then we can only go through each edge once and therefore should pick a simple path. So we can run DFS/BFS from vertex $1$ and take the maximum over all paths to each leaf node. Also $O(N)$.
\subsubsection{Sample Test Cases}
$25;14;21$
\subsubsection{Solve The Full Problem}
The $k\ge2$ case is solved in subtask $1$. So we just need to consider the $k=1$ case.

Again as we have shown in subtask $2$, we just need to find the simple path with the maximum weight sum in the tree. This is known as the (weighted) \textit{diameter}, defined as $$\mathrm{diamater}=\max_{u,v}\mathrm{dist}(u,v),$$ where $u$ and $v$ are vertices in $G$ and $\mathrm{dist}(u,v)$ is the sum of all weights in the simple path between $u$ and $v$ (which is unique as the graph is a tree). Thus we can:
\begin{itemize}
\item Pick any vertex, say $v$.
\item Run DFS/BFS from $v$ to find the farthest vertex $a$.
\item Run DFS/BFS from $a$ to find the farthest vertex $b$. Return the sum of weights of all edges on this path.
\end{itemize}
Note that the second step always finds an endpoint of the diameter, so the third step guarantees the answer. Also note that this still runs in $O(N)$ time.
\newpage
\section*{Appendix: Python Code for Application Question}
\begin{minted}{python}
def solve(n, e, k):
    if k >= 2:
        return sum(w for _, _, w in e)
    g = [[] for _ in range(n + 1)]
    for u, v, w in e:
        g[u].append((v, w))
        g[v].append((u, w))

    def f(start):
        s = [(start, 0, 0)]
        bn = start
        bd = 0
        while s:
            u, p, d = s.pop()
            if d > bd:
                bd = d
                bn = u
            for v, w in g[u]:
                if v != p:
                    s.append((v, u, d + w))
        return bn, bd
    return f(f(1)[0])[1]


n = 6
e = [(1, 2, 3), (2, 3, 1), (1, 4, 7), (1, 5, 9), (5, 6, 5)]
print(solve(n, e, 7))
print(solve(n, e, 1))
\end{minted}
\end{document}
